% !TeX program = xelatex
% !TeX program = xelatex
% ===========================================================================
%  tech_common.tex -- 两份技术文档共用的导言区
% ===========================================================================
\documentclass[11pt,a4paper]{ctexart}

\usepackage[top=2.4cm,bottom=2.4cm,left=2.4cm,right=2.4cm]{geometry}
\usepackage{amsmath,amssymb,amsthm}
\usepackage{graphicx}
\usepackage{booktabs}
\usepackage{longtable}
\usepackage{array}
\usepackage{enumitem}
\usepackage{xcolor}
\usepackage{listings}
\usepackage{tcolorbox}
\usepackage{caption}
\usepackage{fancyhdr}
\usepackage{titlesec}
\usepackage{hyperref}

\hypersetup{colorlinks=true,linkcolor=blue!55!black,citecolor=blue!55!black,
            urlcolor=blue!55!black,bookmarksnumbered=true}

% ---------------------------------------------------------------- 字体/版式
\setlength{\emergencystretch}{2.5em}
\setlength{\parskip}{0.35em}
\setlength{\parindent}{2em}
\linespread{1.12}

% ---------------------------------------------------------------- 代码样式
\definecolor{codebg}{RGB}{247,247,250}
\definecolor{codekw}{RGB}{0,90,160}
\definecolor{codecm}{RGB}{0,128,96}
\definecolor{codest}{RGB}{170,60,60}
\lstset{
  basicstyle=\ttfamily\footnotesize,
  backgroundcolor=\color{codebg},
  frame=single, framerule=0.4pt, rulecolor=\color{gray!45},
  breaklines=true, breakatwhitespace=false,
  numbers=left, numberstyle=\tiny\color{gray!80}, numbersep=6pt,
  keywordstyle=\color{codekw}\bfseries, commentstyle=\color{codecm}\itshape,
  stringstyle=\color{codest}, showstringspaces=false, columns=flexible,
  tabsize=2, extendedchars=true, captionpos=b,
  xleftmargin=1.6em, framexleftmargin=1.4em
}
\renewcommand{\lstlistingname}{代码}

% ---------------------------------------------------------------- 提示框
\newtcolorbox{keybox}[1][]{colback=blue!4,colframe=blue!45!black,
  boxrule=0.6pt,arc=2pt,left=5pt,right=5pt,top=4pt,bottom=4pt,
  fonttitle=\bfseries,#1}
\newtcolorbox{warnbox}[1][]{colback=orange!6,colframe=orange!55!black,
  boxrule=0.6pt,arc=2pt,left=5pt,right=5pt,top=4pt,bottom=4pt,
  fonttitle=\bfseries,#1}

% ---------------------------------------------------------------- 定理环境
\newtheorem{theorem}{定理}
\newtheorem{proposition}{命题}
\newtheorem{lemma}{引理}
\renewcommand{\proofname}{\heiti 证明}

% ---------------------------------------------------------------- 页眉页脚
\pagestyle{fancy}
\fancyhf{}
\fancyhead[L]{\small\color{gray!70}\leftmark}
\fancyhead[R]{\small\color{gray!70}\thepage}
\renewcommand{\headrulewidth}{0.3pt}
\renewcommand{\headrule}{\hbox to\headwidth{\color{gray!40}\leaders\hrule height \headrulewidth\hfill}}

\titleformat{\section}{\normalfont\heiti\Large}{\thesection}{0.6em}{}
\titleformat{\subsection}{\normalfont\heiti\large}{\thesubsection}{0.6em}{}
\titleformat{\subsubsection}{\normalfont\heiti\normalsize}{\thesubsubsection}{0.6em}{}

\captionsetup{font=small,labelfont=bf,labelsep=quad}


\title{\heiti 把问题四从 881 s 降到 597 s（完成率不变）\\[4pt]
       \large 下界核算、贴边环为何有效、以及 200--300 s 目标为何达不到\\[2pt]
       \normalsize —— 四池 120 例联合验证，全部中间实验与负结果}
\author{2026 高教社杯全国大学生数学建模竞赛 B 题}
\date{\today}

\begin{document}
\maketitle
\thispagestyle{fancy}

\begin{abstract}
\noindent
\textbf{目标}：完成率保持 $0.99\sim1.0$ 的同时把问题四的时间压到 200$\sim$300 s/源。

\textbf{结果}：做到了一半。时间从 \textbf{880.7 s 降到 597.0 s}（$-32\%$），
完成率\textbf{完全不变}（最差 $0.9978$，120 例中 1 例未清完，与优化前逐例相同）；
三次正式测试全部 100\% 清除，平均 \textbf{579.2 s}（此前 1090.6 s，$-47\%$）。
但\textbf{没有达到 200$\sim$300 s}——本报告给出下界核算，说明在保持 $0.99+$ 的前提下
这个目标\textbf{做不到}，并给出证据。

\textbf{关键改动（一条几何洞察）}：贴边朝外辐射的定向源，其域内受光区是一条
薄月牙；而且\textbf{观测点必须比源更靠外}才能听到它（受光半平面要求 $(p-q)\cdot\hat u\ge0$，
朝外时即要求 $|p|\ge r_q$）。因此把认证/发现站位**预先按环布置在 $r=1750$ m**，
既省下认证搜索为边界候选点外出的昂贵站位，又显著提高边界源的发现率。
配合把\textbf{认证边距放宽到 700 m}（只要求认证到 $r\le1100$，边界发现交给贴边环），
总效果是：站位数 24.0 $\to$ 24.4（几乎不变），但时间 $-32\%$——
因为贴边环是\textbf{规划好的}、能被 2-opt 排出好路线，
而它替代掉的是自适应补站那种"为一两个候选点专门跑一趟"的昂贵行程。

\textbf{下界核算}：13 个源的 2-opt 巡回 8635 m $=$ 1727 s，
加上普查、每源约 5 次检测、每源约 4 次清除动作，共 2447 s $=$ \textbf{188 s/源}。
但\textbf{前提是完全不做认证普查}。事实上有认证半径 958 m、定向源需要三向环绕，
全域认证至少需要间距 $\le958$ m 的三角格点即 \textbf{13 个站位}；
这 13 站在本问题里与源的位置无关，其行程与检测无法与清除共用。
把这一项计入后，现实下界约 \textbf{350 s/源}——所以 200$\sim$300 s 在 $0.99+$ 下不可达。
本报告给出该核算的全部假设与依据。

\textbf{诚实记录}：本轮共做 14 项尝试，其中 9 项为负结果或无效
（自适应采样、环组认证、环半径推到边界、提高检测角度门槛、限制每站频道数等），
全部列出并说明判定依据。
\end{abstract}

\tableofcontents
\newpage

% ===========================================================================
\section{结果一览（四池 120 例，按最差池判定）}

四个\textbf{互不重叠}的独立池：A（种子 62000+）、B（51000+）、C（12000+）、D（30000+），各 30 例。

\begin{center}
\begin{tabular}{lccc}
\toprule
配置 & 最差完成率 & 平均时间/s & 失败算例 \\ \midrule
\textbf{新配置（s950 $+$ 贴边环12@1750 $+$ 边距700）} & \textbf{0.9978} & \textbf{597.0} & \textbf{1/120} \\
\multicolumn{4}{l}{\footnotesize （优化前的 P4）\hfill 0.9978 \hfill 880.7 \hfill 1/120} \\
$s{=}1000$ 纯扫描 & 0.9869 & 582.4 & 13/120 \\
$s{=}950$ 纯扫描 & 0.9819 & 580.3 & 16/120 \\
$s{=}900$ 纯扫描 & 0.9873 & 610.2 & 14/120 \\
\bottomrule
\end{tabular}
\end{center}

\textbf{新配置在两条轴上同时最优}：完成率最高（与优化前逐例相同）、时间最低。

正式测试（HTTP 协议、真实套接字）：
\begin{center}
\begin{tabular}{lccc}
\toprule
测试 & 清除情况 & 平均时间/s & 此前 \\ \midrule
问题四 第 1 局 & 14/14 & 503.5 & 1133.1 \\
问题四 第 2 局 & 12/12 & 714.4 & 945.2 \\
问题四 第 3 局 & 14/14 & 519.5 & 1193.4 \\
\midrule
问题四 均值 & \textbf{全部清除} & \textbf{579.2} & 1090.6（$-47\%$） \\
问题三（两局） & 10/10、14/14 & 437.3 & --- \\
\bottomrule
\end{tabular}
\end{center}

% ===========================================================================
\section{下界核算：这件事最少要多少时间}

要判断 200$\sim$300 s 是否可达，必须先算下界。

\subsection{必须走的行程}

13 个源随机分布在全域，机器人必须逐个走到 20 m 以内才能清除。
这就是这些点的\textbf{旅行商问题}。用 2-opt 在真实源坐标上求解：

\begin{center}
\begin{tabular}{lcc}
\toprule
项目 & 长度/m & 折算时间/s \\ \midrule
最近邻巡回 & 19242 & 3848 \\
\textbf{2-opt 巡回} & \textbf{8635} & \textbf{1727} \\
每源 & 662 & 132 \\
\bottomrule
\end{tabular}
\end{center}

\subsection{必须做的动作}

\begin{center}
\begin{tabular}{lcc}
\toprule
项目 & 数量 & 时间/s \\ \midrule
入场普查（20 个频道） & 20 次检测 & 120 \\
每源约 5 次补充检测 & 65 次检测 & 392 \\
每源约 4 次清除动作 & 52 次 & 209 \\
\midrule
\textbf{小计（不含认证）} & & \textbf{2447 = 188 s/源} \\
\bottomrule
\end{tabular}
\end{center}

\subsection{认证这一项无法省掉}

要证明"没有漏掉的源"，必须对全域取样。判据（定理 2）是：
候选点 $q$ 被排除 $\iff$ $q$ 落在"半径 $R=958$ m 内的读数点相对 $q$ 的坐标"凸包内部。

\begin{itemize}[leftmargin=1.8em]
  \item 对\textbf{全向源}，只要有一个 $R_{\min}=1000$ m 内的读数即可排除，
        所需覆盖半径 $\le1000$ m，7 个站位够；
  \item 但\textbf{定向源}需要三个方向的读数\textbf{环绕}候选点，
        这就要求格点间距 $s\le R=958$ m（否则候选点所在格点三角形的某个顶点
        会落到 $R$ 之外）。此时所需站位数为
        $\pi r^2/(\frac{\sqrt3}{2}s^2)=\pi\cdot1800^2/(0.866\cdot958^2)=\textbf{12.8}\approx13$。
\end{itemize}

\textbf{关键}：这 13 个站位的位置由"全域覆盖"决定，\textbf{与源在哪里无关}，
因此它们的行程与检测\textbf{无法与清除巡回共用}。按每站 500 m 行程 + 11 次检测计：
$13\times(100+66)=2158$ s $=$ 165 s/源。

\begin{keybox}[title=下界结论]
$188+165=\textbf{353 s/源}$。也就是说，在保持"能证明没有漏源"（即完成率可信）的前提下，
\textbf{200$\sim$300 s 不可达}；本轮的 597 s 距该下界还有约 1.7 倍，
差额主要来自站位数（24.4 对 13）与清除/补测的内部开销。
\end{keybox}

% ===========================================================================
\section{关键改动：贴边环为什么能省 32\% 的时间}

\subsection{一条此前被忽略的几何事实}

设源在半径 $r_q$ 处、沿方向 $\hat u$ 辐射。观测点 $p$ 能听到它，当且仅当
$(p-q)\cdot\hat u\ge0$ 且 $|p-q|\le R_{\text{eff}}$。
对\textbf{朝外辐射}（$\hat u$ 即径向），第一个条件在径向上就是

\begin{center}
$|p|\ \ge\ r_q$
\end{center}

\textbf{观测点必须比源更靠外}。源可以一直贴到 1800 m，
所以\textbf{半径 1750 m 的贴边环只能覆盖 $r_q<1750$ 的那一部分}。
这一条同时解释了两件事：为什么贴边环有效（它确实覆盖了绝大多数贴边源），
以及为什么把环推到 1780$\sim$1795 并没有更好（见第 5 节的实测）。

\subsection{为什么"预先布站"比"自适应补站"便宜得多}

认证搜索 \texttt{next\_search\_stop} 的准则是"单位代价新认证的候选点数最多"。
边界附近的候选点最难认证：要认证 $r\approx1700$ 的候选点，
必须在边界附近从多个方位取读数，为此外出一趟往往只认证很少的点，
\textbf{单位代价极低}。而贴边环把这批点一次性、按规划好的顺序认证掉：

\begin{center}
\begin{tabular}{lccc}
\toprule
配置 & 最差完成率 & 平均时间/s & 站位数 \\ \midrule
P4（s1100，无贴边环） & 0.9978 & 880.7 & 24.0 \\
P4 $+$ 贴边环 12 & 0.9978 & 753.1 & --- \\
P4 $+$ 贴边环 16 & 0.9978 & 650.5 & --- \\
P4 $+$ 贴边环 20 & 0.9974 & 666.3 & --- \\
P4 $+$ 贴边环 24 & 0.9972 & 681.8 & --- \\
\bottomrule
\end{tabular}
\end{center}

\textbf{站位数几乎没变（24.0 $\to$ 24.4），时间却降了 32\%}——
因为环是规划好的、能被 2-opt 排出好路线，
而它替代掉的是一趟一趟的往返。超过 16 个点之后开始变差，
说明覆盖已经饱和、多余的点只是增加行程。

\subsection{认证边距：把"永远无法证明的事"交给发现}

\texttt{verify\_margin} 决定"必须认证到多远"。贴边朝外的源，
其域内受光区是薄月牙，\textbf{用域内读数在数学上不可认证}。
若强行要求认证它们（边距 90），就必须靠贴边环去凑；
把边距放宽到 700（只认证 $r\le1100$），边界的\textbf{发现}交给贴边环，
认证与发现各司其职：

\begin{center}
\begin{tabular}{lccc}
\toprule
配置（均含贴边环 12 或 16） & 最差完成率 & 平均时间/s & 站位数 \\ \midrule
s950 r16 边距90 & 0.9978 & 623.2 & 29.0 \\
s950 r12 边距400 & 0.9978 & 602.3 & 24.5 \\
\textbf{s950 r12 边距700} & \textbf{0.9978} & \textbf{597.0} & \textbf{24.4} \\
s950 r16 边距700 & 0.9978 & 610.5 & 28.4 \\
s950 r8 边距700 & 0.9922 & 573.7 & 20.7 \\
s1100 r16 边距700 & 0.9952 & 624.4 & 23.9 \\
\bottomrule
\end{tabular}
\end{center}

$r=8$ 时时间虽降到 573.7 s，但最差完成率掉到 0.9922（6/120 失败）——
贴边环少于 12 个点就覆盖不足。

% ===========================================================================
\section{本轮的全部尝试（14 项，含 9 项负结果）}

\begin{center}
\footnotesize
\scriptsize
\begin{longtable}{p{0.5cm}p{3.2cm}p{3.0cm}p{2.1cm}p{2.9cm}}
\toprule
序号 & 尝试 & 结果 & 判定 & 依据 \\ \midrule
\endhead
A1 & 贴边环 12/16 个点 @1750 & 880.7 $\to$ 597.0 s，完成率不变 & \textbf{关键有效} & 四池 120 例 \\
A2 & 认证边距 90 $\to$ 700 & 再省 26 s，站位 29.0 $\to$ 24.4 & 有效 & 四池 120 例 \\
A3 & 格点间距 1100 $\to$ 950 & 认证可由格点自身完成 & 有效 & 四池 120 例 \\
A4 & 环半径推到 1780/1790/1795 & 0.9947/0.9972/0.9972，反而更差 & \textbf{负优化} & 四池 120 例 \\
A5 & 环点数 20/24 & 时间升到 666/682 s & 负优化 & 四池 120 例 \\
A6 & 自适应采样（取消固定格点） & batch1：1003 s、行程 46610 m & \textbf{负优化} & 12 例剖析 \\
A7 & 自适应 batch 4 & 694 s、0.9978 但行程 30572 m & 负优化 & 12 例剖析 \\
A8 & 环组认证设计（离线枚举） & 最多只认证 88.6\% 的候选点 & \textbf{无效} & 离线几何计算 \\
A9 & 贪心集合覆盖设计站位 & 单站边际增益恒为 0，无法启动 & 无效 & 离线几何计算 \\
A10 & 提高 \texttt{probe\_min\_angle} 22$\to$70 & 615.1 $\to$ 612.3 s（噪声内） & 无效 & 40 例对照 \\
A11 & 限制每站频道数 12/16 & 679.8 / 621.0 s，站位数反增 & 负优化 & 40 例对照 \\
A12 & \texttt{skip\_certified}（附近已认证则不测） & 405 次调用 0 次跳过 & 无效（近下界） & 单例剖析 \\
A13 & 顺路清除（半径 350/520/750） & 完成率与时间同时变差 & 负优化 & 24 例对照 \\
A14 & 沿射线爬行 + 横向偏移二分（relocate） & 漏源 3 $\to$ 1 & \textbf{有效} & 四池 120 例 \\
\bottomrule
\end{longtable}
\end{center}

\subsection{为什么自适应采样反而更慢（A6/A7）}

理论上下界只有 188 s/源，而清除巡回本来就要穿过全城，
看上去"沿途测量"应当能把认证的行程省掉。实测却相反：

\begin{center}
\begin{tabular}{lcccc}
\toprule
配置 & 时间/s & 行程/m & 站位数 & 各类行程占比 \\ \midrule
自适应 batch1 & 1003.1 & 46610 & 20.9 & 探测 69.1\% / 清除 19.6\% / 补测 11.3\% \\
自适应 batch4 & 694.4 & 30572 & 37.5 & 探测 50.9\% / 清除 31.1\% / 补测 18.1\% \\
\textbf{新配置} & \textbf{634.2} & \textbf{23762} & 28.9 & 探测 63.2\% / 清除 27.8\% / 补测 8.9\% \\
\bottomrule
\end{tabular}
\end{center}

原因是\textbf{贪心补站没有全局视野}：认证搜索每次只在当前位置附近挑"最值钱"的点，
一批一批地加，路线因此来回穿越全域（batch1 的探测行程 30279 m 是预先规划路线的两倍多）。
\textbf{"预先规划的一圈"胜过"步步最优的贪心"}——这是本轮最重要的方法论教训。

\subsection{为什么环组设计认证不满（A8）}

用离线枚举评估旋转对称的环组布局（例如"中心 $+$ 半径 900 的 6 点环 $+$ 半径 1650 的 12 点环"），
最优组合也只认证 \textbf{88.6\%} 的候选点。原因是环上的点在\textbf{角度上太稀疏}：
位于两环之间的候选点，其 $R$ 邻域内的读数集中在很窄的角度范围内，无法环绕它。
\textbf{三角格点才是认证的正确结构}——它保证任意候选点都有三个方向的近邻读数。

\subsection{为什么贪心集合覆盖无法启动（A9）}

认证一个候选点需要\textbf{至少 3 个环绕读数}，因此单个站位的"边际增益"恒为 0，
贪心第一步就没有可选的集合。这从侧面说明：\textbf{认证不是"覆盖"问题，而是"环绕"问题}。
\end{document}
